% ============================================================
%  IACV 2025/26 Project --- Synthetic IPPE-Square Pose Estimation
%  Geometry, two-pose ambiguity, and conditioning experiments
%  Compile with pdfLaTeX. Figures are looked up in figures/ then outputs/.
% ============================================================
\documentclass[11pt,a4paper]{article}

\usepackage[utf8]{inputenc}
\usepackage[T1]{fontenc}
\usepackage[english]{babel}
\usepackage[expansion=false]{microtype}
\usepackage[a4paper,margin=2.6cm]{geometry}
\usepackage{fancyhdr,titlesec,enumitem,parskip}
\setlength{\parskip}{0.5\baselineskip}
\setlength{\headheight}{22pt}
\usepackage{amsmath,amssymb,amsthm,mathtools,bm}
\usepackage{booktabs,array,tabularx,caption,graphicx,placeins}
\usepackage[table]{xcolor}
\definecolor{accent}{HTML}{1F4E79}
\definecolor{accent2}{HTML}{2E75B6}
\definecolor{lightbg}{HTML}{F2F6FA}
\definecolor{greytext}{HTML}{595959}
\usepackage[colorlinks=true,linkcolor=accent,citecolor=accent,urlcolor=accent2]{hyperref}
\usepackage{xurl}
\hypersetup{pdftitle={Synthetic IPPE-Square Pose Estimation},
  pdfsubject={IACV project: planar pose estimation, two-pose ambiguity, and conditioning}}
\urlstyle{tt}

\titleformat{\section}{\normalfont\Large\bfseries\color{accent}}
  {\thesection}{0.6em}{}[{\color{accent}\titlerule[1pt]}]
\titleformat{\subsection}{\normalfont\large\bfseries\color{accent2}}
  {\thesubsection}{0.6em}{}
\titleformat{\subsubsection}{\normalfont\normalsize\bfseries\itshape\color{greytext!90!black}}
  {\thesubsubsection}{0.5em}{}
\titlespacing*{\section}{0pt}{1.6em}{0.8em}
\titlespacing*{\subsection}{0pt}{1.2em}{0.5em}
\titlespacing*{\subsubsection}{0pt}{1em}{0.4em}
\pagestyle{fancy}
\fancyhf{}
\renewcommand{\headrulewidth}{0.4pt}
\renewcommand{\headrule}{\hbox to\headwidth{\color{accent}\leaders\hrule height\headrulewidth\hfill}}
\fancyhead[L]{\small\color{greytext} IACV Project}
\fancyhead[R]{\small\color{greytext} Synthetic IPPE-Square Pose Estimation}
\fancyfoot[C]{\small\thepage}
\theoremstyle{definition}
\newtheorem{definition}{Definition}[section]
\theoremstyle{plain}
\newtheorem{theorem}{Theorem}[section]
\theoremstyle{remark}
\newtheorem*{remark}{Remark}
\newcommand{\code}[1]{\texttt{#1}}
% Long identifiers: \texttt cannot hyphenate, so allow breaks after _ and (
\newcommand{\flag}{\texttt{SOLVEPNP\allowbreak\_IPPE\allowbreak\_SQUARE}}

% --- Notation shortcuts ---
\newcommand{\R}{\mathbb{R}}
\newcommand{\SO}[1]{\mathrm{SO}(#1)}
\newcommand{\SE}[1]{\mathrm{SE}(#1)}
\newcommand{\Rgt}{\mathbf{R}_{\text{gt}}}
\newcommand{\tgt}{\mathbf{t}_{\text{gt}}}
\newcommand{\Pcam}{\mathbf{P}_{\text{cam}}}
\newcommand{\Pmark}{\mathbf{P}_{\text{mk}}}
\newcommand{\KK}{\mathbf{K}}

\newcolumntype{Y}{>{\raggedright\arraybackslash}X}
\newcommand{\thcell}[1]{\textcolor{white}{\textbf{#1}}}
\captionsetup{font=small,labelfont=bf}
\setlist[itemize]{itemsep=2pt,topsep=3pt}
\setlist[enumerate]{itemsep=2pt,topsep=3pt}
\setlength{\emergencystretch}{2em}

% Arguments: filename, description of the expected plot, caption, label.
% The figure is searched in figures/ and then in outputs/, which is where
% experiment1.py writes it and which is not tracked by Git.
\newcommand{\resultfigure}[4]{%
  \begin{figure}[h!]
  \centering
  \IfFileExists{figures/#1}{%
    \includegraphics[width=0.88\linewidth,height=0.6\textheight,keepaspectratio]{figures/#1}%
  }{%
    \IfFileExists{outputs/#1}{%
      \includegraphics[width=0.88\linewidth,height=0.6\textheight,keepaspectratio]{outputs/#1}%
    }{%
      \fcolorbox{accent2}{lightbg}{%
        \begin{minipage}{0.89\linewidth}
        \centering\vspace{0.7cm}
        {\color{accent}\bfseries Figure pending}\par\medskip
        #2\par\medskip
        {\footnotesize Expected file: \texttt{\detokenize{#1}}}\par\smallskip
        {\footnotesize Regenerate it by running the corresponding experiment script.}
        \vspace{0.7cm}
        \end{minipage}}
    }%
  }
  \caption{#3}\label{#4}
  \end{figure}%
}

\begin{document}
\begin{titlepage}
\centering
\vspace*{3cm}
{\color{accent2}\bfseries\large IACV 2025/26 PROJECT}\par
\vspace{0.6cm}
{\color{accent}\bfseries\Huge Synthetic IPPE-Square}\par
\vspace{0.3cm}
{\color{accent}\bfseries\Huge Pose Estimation}\par
\vspace{0.8cm}
{\color{greytext}\itshape\Large Technical Project Report}\par
\vspace{0.5cm}
{\color{greytext}\large Planar geometry, the two-pose ambiguity, and conditioning}\par
\vspace{0.4cm}
{\color{greytext} Mathematical foundations, implementation, and synthetic experiments}
\vfill
{\color{greytext}\large Author names\par}
\vfill
{\color{greytext} September 2026}
\end{titlepage}

\tableofcontents
\thispagestyle{empty}
\clearpage

% ============================================================
\section{Overview and Notation}
\label{sec:overview}
% ============================================================

We synthesize the image a calibrated pinhole camera would observe of a
known planar square marker placed at a known pose. There is no detection or
image-processing stage: all 3D--2D corner correspondences are generated
mathematically.

\begin{enumerate}[label=\textbf{Step \arabic*.}, leftmargin=*]
  \item Generate a 3D plane (the marker plane).
  \item Define a square on that plane of side $L$.
  \item Define a camera with known intrinsics $\KK$.
  \item Choose a known marker pose $(\Rgt, \tgt)$ relative to the camera.
  \item Project the four 3D corners into the image.
  \item Recover both planar IPPE-square pose candidates from the correspondences.
  \item Compare the custom primary candidate to OpenCV's
        \flag{} in the noiseless case.
\end{enumerate}

Section~\ref{sec:pipeline} follows these seven steps in order. Two experiments
then build on the finished pipeline: Section~\ref{sec:experiment1} sweeps the
camera viewpoint and studies the two planar candidates, and
Section~\ref{sec:experiment2} analyses the reprojection Jacobian and compares
its uncertainty prediction with a Monte Carlo simulation.

\subsection{Notation}
\begin{itemize}[leftmargin=*]
  \item Vectors in bold lowercase: $\mathbf{t}$. Matrices in bold uppercase: $\mathbf{R}, \KK$.
  \item Points carry a subscript naming their frame:
        $\Pmark$ is in the marker frame, $\Pcam$ in the camera frame.
  \item Pixel coordinates: $\mathbf{p} = (u, v)^\top$.
  \item Coordinate frames are right-handed: in the camera frame, $X$ points right,
        $Y$ points down, $Z$ points into the scene.
\end{itemize}

% ============================================================
\section{The Synthetic Pipeline}
\label{sec:pipeline}
% ============================================================

Steps 1--5 build a scene and the image a calibrated camera would record of it;
steps 6--7 recover the pose from that image and check the result. Each
subsection below corresponds to one step of the list in
Section~\ref{sec:overview} and to one stage of \code{pipeline.py}, which prints
every intermediate quantity when run.

% ------------------------------------------------------------
\subsection{Step 1 --- The 3D Marker Plane}
% ------------------------------------------------------------

The marker is a flat object, so the four corners are coplanar. We attach a
coordinate frame to the marker (the \emph{marker frame}) with the plane
chosen as $Z = 0$:
\begin{equation}
  \text{marker plane} \;=\; \{\, (X, Y, 0) \,:\, X, Y \in \R \,\}.
\end{equation}
The origin is placed at the geometric center of the future square. The
$X$ axis points right, $Y$ points up (in the plane), and $Z$ is the outward
normal to the marker surface.

\paragraph{Why $Z = 0$ matters.}
A general projection of a 3D point requires the full $3\times 3$ rotation
matrix. If the third coordinate is identically zero, the third column of
$\mathbf{R}$ multiplies $0$ and drops out:
\begin{equation}
  \mathbf{R} \begin{pmatrix} X \\ Y \\ 0 \end{pmatrix}
  \;=\; X \mathbf{r}_1 + Y \mathbf{r}_2,
\end{equation}
where $\mathbf{r}_1, \mathbf{r}_2$ are the first two columns of $\mathbf{R}$. This
collapse is the structural reason a single $3\times 3$ matrix --- a
\emph{homography} --- captures all the geometry of a planar marker, which
is what later pose-estimation steps will exploit.

% ------------------------------------------------------------
\subsection{Step 2 --- The Square on the Plane}
% ============================================================

On the plane $Z = 0$ we place a square of side $L$, centered at the marker
origin and axis-aligned to the marker frame. The four corner coordinates,
in the order used by OpenCV's \flag, are
\begin{equation}
  \Pmark^{(0)} = \begin{pmatrix} -L/2 \\ +L/2 \\ 0 \end{pmatrix}, \;
  \Pmark^{(1)} = \begin{pmatrix} +L/2 \\ +L/2 \\ 0 \end{pmatrix}, \;
  \Pmark^{(2)} = \begin{pmatrix} +L/2 \\ -L/2 \\ 0 \end{pmatrix}, \;
  \Pmark^{(3)} = \begin{pmatrix} -L/2 \\ -L/2 \\ 0 \end{pmatrix}.
\end{equation}
The parenthesized superscript is the corner index; the subscript ``mk''
indicates the marker frame.

\paragraph{Corner ordering matters.}
The order of the 3D list must match the order of the 2D observations,
otherwise the correspondences are scrambled and the algorithm finds the
wrong pose (or fails). The order above (clockwise starting at top-left, when
the marker frame is drawn with $X$ right and $Y$ up and seen from $+Z$) is the
convention shared by OpenCV's ArUco detector and \flag.
The sign of the shoelace sum over the four corners is negative, which is the
quick numerical check for this orientation.

% ------------------------------------------------------------
\subsection{Step 3 --- The Pinhole Camera}
% ============================================================

\subsubsection{Geometric model}

A pinhole camera maps points in the camera frame to points on a 2D image.
The camera origin is at the world origin; the image plane sits at $Z = f$,
where $f$ is the focal length. By similar triangles, a 3D point
$\Pcam = (X, Y, Z)^\top$ projects to
\begin{equation}
  x = f\cdot\frac{X}{Z}, \qquad y = f\cdot\frac{Y}{Z}.
  \label{eq:similar-triangles}
\end{equation}
The division by $Z$ is the \emph{perspective divide}: it is the only
non-linearity in the model and is responsible for foreshortening (farther
objects look smaller).

\subsubsection{The intrinsic matrix \texorpdfstring{$\KK$}{K}}

Real images are measured in pixels. The focal length is encoded as
$f_x, f_y$ (in pixels, possibly different along the two axes if the pixel
elements are not square), and the optical axis hits the image at the
\emph{principal point} $(c_x, c_y)$. The pixel coordinates are
\begin{equation}
  u = f_x \cdot \frac{X}{Z} + c_x, \qquad
  v = f_y \cdot \frac{Y}{Z} + c_y.
  \label{eq:pixel-projection}
\end{equation}
Collected into a matrix:
\begin{equation}
  \boxed{\quad
    \KK \;=\; \begin{pmatrix} f_x & 0 & c_x \\ 0 & f_y & c_y \\ 0 & 0 & 1 \end{pmatrix}
  \quad}
\end{equation}
so that the projection has the compact homogeneous form
\begin{equation}
  \tilde{\mathbf{p}} \;=\; \KK\,\Pcam,
  \qquad
  \mathbf{p} \;=\; \begin{pmatrix} \tilde p_1/\tilde p_3 \\ \tilde p_2/\tilde p_3 \end{pmatrix}.
\end{equation}
The first is a linear matrix product, the second is the perspective divide.

\subsubsection{Normalized image coordinates}

Apply $\KK^{-1}$ to a pixel $\mathbf{p}$:
\begin{equation}
  \begin{pmatrix} x_n \\ y_n \\ 1 \end{pmatrix}
  \;=\; \KK^{-1}
  \begin{pmatrix} u \\ v \\ 1 \end{pmatrix}
  \;=\;
  \begin{pmatrix} (u - c_x)/f_x \\ (v - c_y)/f_y \\ 1 \end{pmatrix}.
\end{equation}
The pair $(x_n, y_n)$ are the \emph{normalized image coordinates}: the
location where the ray through pixel $\mathbf{p}$ crosses the plane
$Z = 1$ in camera space. They are dimensionless (no pixels, no focal
length) and encode pure ray directions. Pose-estimation algorithms
work in this representation so that the camera-specific quantities $\KK$
do not contaminate the geometric reasoning.

% ------------------------------------------------------------
\subsection{Step 4 --- Ground-Truth Pose}
% ============================================================

A point in the marker frame is brought into the camera frame by a
rigid-body transform:
\begin{equation}
  \boxed{\quad \Pcam \;=\; \mathbf{R}\,\Pmark + \mathbf{t} \quad}
  \label{eq:rigid-transform}
\end{equation}
with $\mathbf{R} \in \SO{3}$ and $\mathbf{t} \in \R^3$. The pair
$(\mathbf{R}, \mathbf{t})$ is an element of the Special Euclidean group
$\SE{3}$, and is what we mean by \emph{pose}.

\subsubsection{The rotation group \texorpdfstring{$\SO{3}$}{SO(3)}}

A $3\times 3$ matrix $\mathbf{R}$ is a rotation if and only if
\begin{equation}
  \mathbf{R}^\top \mathbf{R} \;=\; \mathbf{I},
  \qquad
  \det(\mathbf{R}) \;=\; +1.
  \label{eq:so3-constraints}
\end{equation}
The first condition (orthonormal columns) preserves lengths and angles; the
second excludes reflections, which would flip a right-handed frame into a
left-handed one.

\paragraph{Degrees of freedom.}
A $3\times 3$ matrix has $9$ entries. The orthogonality condition
$\mathbf{R}^\top \mathbf{R} = \mathbf{I}$ imposes $6$ independent constraints. That
leaves \textbf{3 DOF} --- the minimum number of parameters needed to
describe a rotation in 3D.

\subsubsection{Representation 1: Euler angles}

Decompose any rotation as three elementary rotations around the coordinate
axes. The code uses SciPy's lower-case \emph{extrinsic xyz} convention: the
rotations are applied about fixed camera axes, and the resulting active
matrix is
\begin{equation}
  \mathbf{R}(\alpha,\beta,\gamma) \;=\; \mathbf{R}_z(\gamma)\,\mathbf{R}_y(\beta)\,\mathbf{R}_x(\alpha).
\end{equation}
where each $\mathbf{R}_\bullet(\cdot)$ is a $3\times 3$ planar rotation, e.g.
\begin{equation}
  \mathbf{R}_x(\alpha) =
  \begin{pmatrix} 1 & 0 & 0 \\ 0 & \cos\alpha & -\sin\alpha \\ 0 & \sin\alpha & \cos\alpha \end{pmatrix}.
\end{equation}
\textbf{Pros:} human-intuitive --- $(15^\circ, 10^\circ, 5^\circ)$ is easy to picture.\\
\textbf{Cons:} suffer from \emph{gimbal lock} (a configuration where two
axes align and one DOF is lost) and from non-unique decompositions. We use
Euler angles only to \emph{specify} the ground-truth pose conveniently; we
never use them for optimization.

\subsubsection{Facing the camera}

One detail is easy to miss. If $\mathbf{R} = \mathbf{I}$, the marker axes
coincide with the camera axes, so the marker normal $+Z$ points \emph{into}
the scene, away from the camera: we would be looking at the back of the
marker. A marker that faces the camera, as a detected ArUco marker always
does, is therefore near a $180^\circ$ rotation, not near the identity. We
build the pose as
\begin{equation}
  \Rgt \;=\; \mathbf{R}(\alpha,\beta,\gamma)\, \mathbf{F},
  \qquad
  \mathbf{F} = \operatorname{diag}(1,-1,-1),
  \label{eq:face-flip}
\end{equation}
where $\mathbf{F}$ is a $180^\circ$ rotation about $X$ (a proper rotation:
$\det \mathbf{F} = +1$). Because $\mathbf{F}$ is applied first, in the marker
frame, the Euler angles keep their intended meaning: a tilt away from the
fronto-parallel view.

\subsubsection{Our concrete ground-truth pose}
\label{sec:gt-pose}

In the code we choose
\begin{align}
  \text{Euler (xyz, deg)} & : \quad (\alpha, \beta, \gamma) = (15,\,10,\,5) \\
  \tgt & : \quad (0.05,\, 0.02,\, 0.50)^\top \;\text{m}
\end{align}
which, through \eqref{eq:face-flip}, yields
\[
  \Rgt \approx
  \begin{pmatrix}
    +0.9811 & +0.0394 & -0.1897 \\
    +0.0858 & -0.9662 & +0.2432 \\
    -0.1736 & -0.2549 & -0.9513
  \end{pmatrix}.
\]
The Euler part alone rotates by $\approx 18.33^\circ$, which is \emph{not}
$15^\circ + 10^\circ + 5^\circ$ because Euler rotations do not compose by
simple addition. The third column of $\Rgt$ is the marker normal expressed in
the camera frame, $(-0.190,\, 0.243,\, -0.951)^\top$; its negative $Z$
component confirms that the normal points back towards the camera.

\paragraph{Viewing angle.} Throughout this report the \emph{viewing angle} is
the angle between the marker normal and the line from the marker centre to the
camera centre: $0^\circ$ is fronto-parallel, and a value above $90^\circ$
would mean we are looking at the back of the marker. For the pose above it is
$17.26^\circ$, so the baseline scene is a mildly oblique but ordinary view,
directly comparable with the sweep of Section~\ref{sec:experiment1}.

% ------------------------------------------------------------
\subsection{Step 5 --- Forward Projection}
% ============================================================

Steps 1--4 fix all the inputs. Step 5 chains them into the observed image
points. For each marker corner $\Pmark^{(i)}$:
\begin{equation}
  \boxed{\quad
    \mathbf{p}^{(i)}
    \;=\;
    \pi\!\Big(\, \KK\,(\,\Rgt \Pmark^{(i)} + \tgt\,) \,\Big)
  \quad}
  \label{eq:full-projection}
\end{equation}
where the projection $\pi:\R^3 \to \R^2$ is the perspective divide
\begin{equation}
  \pi\!\begin{pmatrix} \tilde x \\ \tilde y \\ \tilde z \end{pmatrix}
  \;=\;
  \begin{pmatrix} \tilde x / \tilde z \\ \tilde y / \tilde z \end{pmatrix}.
\end{equation}

\paragraph{In the code, Step 5 is split in two sub-steps:}
\begin{enumerate}[label=(\alph*)]
  \item \textbf{Rigid transform.} $\Pcam^{(i)} \;=\; \Rgt \Pmark^{(i)} + \tgt$.
        Brings each marker corner from the marker frame into the camera frame.
  \item \textbf{Pinhole projection.} $\mathbf{p}^{(i)} \;=\; \pi(\KK\,\Pcam^{(i)})$.
        Maps the camera-frame point to pixel coordinates via
        \eqref{eq:pixel-projection}.
\end{enumerate}

\paragraph{Possible sanity checks at this point:}
\begin{itemize}[leftmargin=*]
  \item All $\Pcam^{(i)}$ have $Z > 0$ (in front of the camera).
  \item All $\mathbf{p}^{(i)}$ fall inside $[0, W) \times [0, H)$ (visible in the image, W as width, H as height).
  \item The four $\mathbf{p}^{(i)}$ form a \emph{convex quadrilateral} that looks
        like a perspective-warped square.
\end{itemize}

% ------------------------------------------------------------
\subsection{Step 6 --- IPPE-Square Pose Recovery}
\label{sec:ippe}
% ------------------------------------------------------------

The marker is planar, so its normalized-image projection is a homography
\begin{equation}
  \begin{pmatrix} x_n \\ y_n \\ 1 \end{pmatrix}
  \sim \mathbf{H}\begin{pmatrix} X \\ Y \\ 1 \end{pmatrix}.
\end{equation}
Our solver is written from scratch in \code{ippe\_square.py} and never calls
OpenCV. It follows the planar-homography formulation of IPPE, and this section
describes what the implementation actually computes, function by function.

\subsubsection{Input validation}

\code{IPPE\_SQUARE} is not a generic coplanar solver: it assumes the four
object points are the centred square of Step 2, in that exact order. Feeding it
a different layout still produces a pose, only a meaningless one. The solver
therefore first checks that the points are finite, lie on $Z=0$, and match the
canonical square of side $L = \|\Pmark^{(1)} - \Pmark^{(0)}\|$, and refuses
anything else.

\subsubsection{From pixels to a homography}

Pixels are first mapped to normalized coordinates with $\KK^{-1}$, so the
intrinsics disappear from everything that follows. The homography is then
estimated by the standard DLT: each correspondence $\big((X,Y),(x_n,y_n)\big)$
contributes two rows to a $8\times9$ design matrix, and $\mathbf{H}$ is read off
the right singular vector of its smallest singular value. The homogeneous
scale is fixed by dividing by $h_{33}$, which is safe here because the square is
centred at the origin and lies in front of the camera.

\subsubsection{The two quantities IPPE needs}

IPPE is a \emph{local} method: it uses the homography only through its value
and its first derivative at one point, here the square centre $(X,Y)=(0,0)$.
Since the centre is the origin, its normalized projection is read directly
off the last column,
\begin{equation}
  \mathbf{v} = (h_{13},\, h_{23})^\top ,
\end{equation}
and the derivative follows from the quotient rule applied to
$x_n(X,Y) = (h_{11}X + h_{12}Y + h_{13})/(h_{31}X + h_{32}Y + 1)$, evaluated at
the origin:
\begin{equation}
  \mathbf{J}_{\mathbf H} =
  \begin{pmatrix}
    h_{11} - h_{13}h_{31} & h_{12} - h_{13}h_{32} \\
    h_{21} - h_{23}h_{31} & h_{22} - h_{23}h_{32}
  \end{pmatrix}.
  \label{eq:homography-jacobian}
\end{equation}

\subsubsection{The IPPE construction}

Let $\mathbf{R}_{\mathbf v}$ be the smallest rotation taking $\mathbf{e}_3$ onto
$(\mathbf{v},1)^\top$, computed in closed form by Rodrigues' formula. Working in
the frame it defines removes the effect of the marker being off the optical
axis. Writing
\begin{equation}
  \mathbf{B} =
  \begin{pmatrix} 1 & 0 & -v_1 \\ 0 & 1 & -v_2 \end{pmatrix}
  \big[\mathbf{R}_{\mathbf v}\big]_{:,1:2},
  \qquad
  \mathbf{A} = \mathbf{B}^{-1}\mathbf{J}_{\mathbf H},
\end{equation}
the matrix $\mathbf{A}$ is, up to scale, the upper-left $2\times2$ block of the
rotation we are looking for. Its largest singular value $\gamma$ supplies that
scale, and we set $\mathbf{R}_{2\times2} = \mathbf{A}/\gamma$.

The rotation is then completed column by column. Orthonormality forces
\begin{equation}
  \mathbf{b}\,\mathbf{b}^\top
  = \mathbf{I} - \mathbf{R}_{2\times2}^\top\mathbf{R}_{2\times2},
  \label{eq:rank-one}
\end{equation}
whose right-hand side is positive semidefinite of rank one; $\mathbf{b}$ is
recovered from its largest eigenpair. The remaining entries come from the
cross product of the first two columns, which must be the third. Both signs of
$\mathbf{b}$ satisfy \eqref{eq:rank-one}, and this is precisely where the
planar ambiguity enters: the two candidates are
\begin{equation}
  \mathbf{R}^{\pm} = \mathbf{R}_{\mathbf v}
  \begin{pmatrix}
    \mathbf{R}_{2\times2} & \pm\mathbf{c} \\[2pt]
    \pm\mathbf{b}^\top & a
  \end{pmatrix},
\end{equation}
with $(\mathbf{c},a)$ obtained from that cross product. At a fronto-parallel
view the right-hand side of \eqref{eq:rank-one} vanishes, so $\mathbf{b} =
\mathbf{0}$ and the two candidates collapse into one; Section~\ref{sec:experiment1}
returns to this.

\subsubsection{Translation, ranking, and the decision not to refine}

With a rotation fixed, translation is linear. Each corner requires its camera
ray $\mathbf{q}^{(i)} = (x_n^{(i)}, y_n^{(i)}, 1)^\top$ to be parallel to the
transformed point, that is
\begin{equation}
  \mathbf{q}^{(i)} \times \big(\mathbf{R}\,\Pmark^{(i)} + \mathbf{t}\big)
  = \mathbf{0}
  \quad\Longleftrightarrow\quad
  [\mathbf{q}^{(i)}]_\times\,\mathbf{t}
  = -[\mathbf{q}^{(i)}]_\times\,\mathbf{R}\,\Pmark^{(i)} .
  \label{eq:translation}
\end{equation}
Only two of the three rows per corner are independent: writing
$\mathbf{P}=\mathbf{R}\,\Pmark^{(i)}+\mathbf{t}$, the first two rows read
$P_2 - y_n P_3 = 0$ and $P_1 - x_n P_3 = 0$, while the third,
$x_n P_2 - y_n P_1 = 0$, is $x_n$ times the first minus $y_n$ times the second. We keep the first two, so stacking
the four corners gives an $8\times3$ system, solved in the least-squares sense.
This is the same system OpenCV solves; Section~\ref{sec:step7} explains why the
choice matters.

Both candidates are then reprojected, ranked by pixel RMSE, and returned
together. We deliberately apply \emph{no} nonlinear refinement: an
unconstrained optimization started from the two candidates can pull both to the
same pose, which would erase the very ambiguity Experiment 1 sets out to
measure.

\subsubsection{Reference implementation}

The result is compared, in the same noiseless scene, with OpenCV's
\code{solvePnPGeneric} called with the \flag{} flag, using the same canonical
point order. OpenCV serves only as an external reference: no part of the solver
depends on it.

% ------------------------------------------------------------
\subsection{Step 7 --- Validation Against OpenCV and the Ground Truth}
\label{sec:step7}
% ------------------------------------------------------------

Table~\ref{tab:baseline} collects the four poses returned for the pose of
Section~\ref{sec:gt-pose}. The primary candidate reproduces the ground truth to machine
precision, and the per-corner residuals are of order $10^{-13}$ px, so the
solver is exact in the noiseless case rather than merely accurate.

\begin{table}[h!]
\centering\small
\begin{tabularx}{\linewidth}{@{}YYrrr@{}}
\toprule
\rowcolor{accent}\thcell{Candidate}&\thcell{Solver}&\thcell{RMSE (px)}&\thcell{Rot.\ error ($^\circ$)}&\thcell{Transl.\ error (mm)}\\
\midrule
1&custom&$6.4\cdot10^{-14}$&$0$&$1.1\cdot10^{-13}$\\
\rowcolor{lightbg} 2&custom&$3.33825$&$34.521$&$3.101$\\
1&OpenCV&$2.2\cdot10^{-14}$&$0$&$1.1\cdot10^{-13}$\\
\rowcolor{lightbg} 2&OpenCV&$3.33825$&$34.521$&$3.101$\\
\bottomrule
\end{tabularx}
\caption{Baseline scene: both IPPE-square candidates from the custom solver and
from OpenCV, each compared with the ground-truth pose.}
\label{tab:baseline}
\end{table}

The candidate-1 values should not be compared with each other. The exact error
is zero, since the pose fits the data perfectly, so what is printed is pure
floating-point rounding. A double carries about $16$ significant digits, and
with pixel coordinates near $500$ px the smallest representable step is about
$500 \times 2.2\cdot10^{-16} \approx 10^{-13}$ px. Both $6.4\cdot10^{-14}$ and
$2.2\cdot10^{-14}$ lie below that floor, so both are zero for every practical
purpose; which one comes out larger depends only on the order of the
floating-point operations inside each solver. The same absolute difference,
about $2\cdot10^{-14}$ px, separates the two candidate-2 values, where it is
simply invisible next to $3.3$ px.

Two observations are worth making. First, the two solvers agree on \emph{both}
candidates: to $10^{-13}$ mm on the primary pose, and to every printed digit on
the mirror one, which validates the custom implementation. Second, the
mirror candidate is $34.52^\circ$ away from the ground truth, that is exactly
twice the $17.26^\circ$ viewing angle of the scene; Section~\ref{sec:experiment1}
shows this is not a coincidence.

% \paragraph{How the two solvers were brought into agreement.} A first version of our solver kept all three cross-product rows of \eqref{eq:translation}, a $12\times3$ system, and agreed with OpenCV on the primary pose but not on the mirror one ($3.33836$ against $3.33825$ px, translations $3.074$ against $3.101$ mm from the ground truth). We located the cause by elimination. It was not the error metric: recomputing OpenCV's pose with our formula reproduced its number to every digit. It was not the homography: OpenCV builds it in closed form for the square, we use a DLT, and the two agree to $4\cdot10^{-14}$. It was not $\mathbf{b}$ either, which agreed to $2\cdot10^{-15}$. It was the translation. The third row is a linear combination of the other two, so it adds no information, but in a least-squares fit it still reweights the residuals. When the rotation fits the data exactly the two systems share a solution --- hence the agreement on the primary pose --- but the mirror rotation never fits exactly, and there they gave translations $2.8\cdot10^{-5}$ m apart. Neither version is more correct, since each minimizes a different algebraic residual; keeping only the first two rows, as OpenCV does, makes the two implementations agree on both candidates, which is the version reported in Table~\ref{tab:baseline}.

\paragraph{A limitation of OpenCV's IPPE at $180^\circ$.} While comparing the
two solvers along the camera arc of Section~\ref{sec:experiment1}, OpenCV
returned wrong poses: about $113$ px of reprojection error on data our solver
fits to $10^{-13}$ px. The data were not at fault, since OpenCV's own
\code{SOLVEPNP\_SQPNP} solves them to $10^{-14}$ px, and the generic
\code{SOLVEPNP\_IPPE} flag fails in the same way as \flag. The translation it
returned was correct and only the rotation was wrong, which points at the
output stage rather than at the IPPE construction. The culprit is the private
routine that converts the rotation matrix into a rotation vector. It uses the
textbook formula
\begin{equation}
  \omega = \arccos\!\Big(\tfrac{\operatorname{tr}\mathbf{R}-1}{2}\Big),
  \qquad
  \boldsymbol{\omega}
  = \frac{\omega}{2\sin\omega}
    \begin{pmatrix}
      R_{32}-R_{23}\\ R_{13}-R_{31}\\ R_{21}-R_{12}
    \end{pmatrix},
\end{equation}
with no special case near $\omega=180^\circ$, where both $\sin\omega$ and the
antisymmetric part of $\mathbf{R}$ vanish and the axis becomes $0/0$. A marker
facing the camera is close to a half turn, and every pose on our arc is
\emph{exactly} one, so it is lost entirely; OpenCV's own \code{cv2.Rodrigues},
which handles this case, recovers the same rotations without error. The
degradation is gradual as the angle approaches $180^\circ$:

\begin{table}[h!]
\centering\small
\begin{tabularx}{\linewidth}{@{}Yrr@{}}
\toprule
\rowcolor{accent}\thcell{Euler angles in the configuration}&\thcell{Rotation angle}&\thcell{OpenCV best RMSE (px)}\\
\midrule
$(15^\circ,10^\circ,5^\circ)$, the baseline & $165.5^\circ$ & $2\cdot10^{-14}$\\
\rowcolor{lightbg} $(5^\circ,3^\circ,1^\circ)$ & $175.0^\circ$ & $4\cdot10^{-13}$\\
$(1^\circ,0^\circ,0^\circ)$ & $179.0^\circ$ & $6\cdot10^{-12}$\\
\rowcolor{lightbg} $(0^\circ,0^\circ,0^\circ)$ & $180.0^\circ$ & $1.21$\\
\bottomrule
\end{tabularx}
\caption{OpenCV's IPPE-square accuracy as the ground-truth rotation approaches
a half turn. With the facing flip of \eqref{eq:face-flip}, zero Euler angles
give exactly $180^\circ$.}
\label{tab:opencv180}
\end{table}

Two consequences follow. The baseline, at $165.5^\circ$, is unaffected, so the
validation above stands. But the natural test of setting all three Euler angles
to zero would make OpenCV fail while our solver remains exact, and it is one
more reason Experiment 1 relies on our implementation: with OpenCV's IPPE the
sweep could not have been run at all.

% ============================================================
\section{Experiment 1 --- The Two-Pose Ambiguity}
\label{sec:experiment1}
% ============================================================

A planar marker admits two poses that explain the same image almost equally
well. This experiment measures \emph{how} equally, as a function of the viewing
angle: when the two candidates are easy to tell apart, and when they are not.

% ------------------------------------------------------------
\subsection{The quantity we plot}
\label{sec:rmse}
% ------------------------------------------------------------

Everything below is measured with the reprojection error, so it is worth
stating exactly what that number is. Given an estimated pose
$(\mathbf{R},\mathbf{t})$, we project the four known corners again and compare
them with the observed pixels. Each corner contributes a two-component
residual,
\begin{equation}
  \mathbf{e}^{(i)}
  = \underbrace{\pi\!\big(\KK(\mathbf{R}\,\Pmark^{(i)} + \mathbf{t})\big)}
    _{\text{where the pose says corner } i \text{ should be}}
  \;-\;
  \underbrace{\mathbf{p}^{(i)}}_{\text{where it is observed}}
  \;=\;
  \begin{pmatrix} \Delta u^{(i)} \\ \Delta v^{(i)} \end{pmatrix},
  \label{eq:residual}
\end{equation}
so four corners give $2N = 8$ scalar residuals in total. Summing their squares
gives the reprojection cost itself,
\begin{equation}
  E(\mathbf{R},\mathbf{t}) = \sum_{i=1}^{N}\|\mathbf{e}^{(i)}\|^2
  = \sum_{i=1}^{N}\Big[(\Delta u^{(i)})^2+(\Delta v^{(i)})^2\Big],
  \label{eq:cost}
\end{equation}
which is the quantity Section~\ref{sec:experiment2} minimizes. What we plot is
its root mean square,
\begin{equation}
  \boxed{\quad
  \mathrm{RMSE}
  =\sqrt{\frac{E(\mathbf{R},\mathbf{t})}{2N}}
  =\sqrt{\frac{1}{2N}\sum_{i=1}^{N}
    \Big[(\Delta u^{(i)})^2+(\Delta v^{(i)})^2\Big]}\,,
  \qquad N = 4 .
  \quad}
  \label{eq:rmse}
\end{equation}

\paragraph{Why not just $E$?} The two carry the same information, since
$\mathrm{RMSE}=\sqrt{E/2N}$ is a monotone function of $E$, but they are not
equally readable. The cost $E$ is measured in squared pixels and grows with the
number of corners, so a value of $89$ says nothing on its own. Taking the mean
and then the square root undoes both effects: the result is in \emph{pixels},
it does not depend on how many points were used, and it answers a concrete
question --- \emph{if this pose were the true one, by how many pixels would a
corner coordinate be misplaced?} That also makes it directly comparable with
the corner noise $\sigma_{px}$, which is the comparison the discussion below
rests on. A perfect pose gives zero either way.

\begin{remark}
The normalization is by $2N$, the number of scalar residuals, not by $N$, the
number of corners. This is the convention OpenCV reports, so our numbers and
its own are directly comparable; an RMSE of Euclidean per-corner distances
would be a factor $\sqrt2$ larger.
\end{remark}

% ------------------------------------------------------------
\subsection{Setup}
% ------------------------------------------------------------

The camera centre moves on an arc of fixed radius $r = 0.5$ m around the
marker, always aimed at the marker centre, sampling one viewpoint per degree
from $0^\circ$ (fronto-parallel) to $89^\circ$ (almost edge-on). No noise is
added: the image points are the exact projections. At each viewpoint we solve
with our own IPPE implementation of Section~\ref{sec:ippe}, keep \emph{both}
candidates, and record the reprojection RMSE and the pose error of each. Using
our solver rather than OpenCV is what makes the experiment possible at all:
both candidates are returned unrefined, so neither is silently discarded.

\paragraph{One arc, not a hemisphere.} The camera path has a single free
parameter. Both the azimuth, fixed at $0^\circ$, and the radius, fixed at
$0.5$ m, are held constant, so the camera slides along one meridian rather than
covering the whole front hemisphere. The square is symmetric under
$90^\circ$ in-plane rotations, so other azimuths mainly relabel the corners and
give equivalent curves; the radius, by contrast, changes the apparent size of
the marker and therefore the scale of the reprojection errors. The results
below should be read as belonging to this configuration.

% ------------------------------------------------------------
\subsection{Results}
% ------------------------------------------------------------

Figure~\ref{fig:sweep} shows the two curves. The true pose (candidate 1) sits
at $\approx 10^{-13}$ px everywhere, as it must without noise. The mirror pose
(candidate 2) starts at zero and grows monotonically to $12.15$ px.

\resultfigure{exp1_sweep.png}
 {Reprojection RMSE of the true and mirror candidates against the viewing angle.}
 {Reprojection RMSE of the two IPPE-square candidates against the viewing
 angle, one sample per degree.}
 {fig:sweep}

\begin{table}[h!]
\centering\small
\begin{tabularx}{\linewidth}{@{}Yrr@{}}
\toprule
\rowcolor{accent}\thcell{Viewing angle ($^\circ$)}&\thcell{Mirror RMSE (px)}&\thcell{Mirror rot.\ error ($^\circ$)}\\
\midrule
0&$0.000$&$0$\\
\rowcolor{lightbg} 1&$0.197$&$2$\\
5&$0.986$&$10$\\
\rowcolor{lightbg} 15&$2.939$&$30$\\
30&$5.736$&$60$\\
\rowcolor{lightbg} 45&$8.240$&$90$\\
60&$10.281$&$120$\\
\rowcolor{lightbg} 75&$11.659$&$150$\\
89&$12.153$&$178$\\
\bottomrule
\end{tabularx}
\caption{How the mirror candidate separates from the true pose, in the image
and in pose space. The rotation error is exactly twice the viewing angle at
every sampled point.}
\label{tab:sweep}
\end{table}

% ------------------------------------------------------------
\subsection{Discussion}
% ------------------------------------------------------------

\subsubsection{Why the two poses coincide at $0^\circ$}

The collapse is visible directly in the algebra of Section~\ref{sec:ippe}. The
two candidates differ only through the vector $\mathbf{b}$, recovered from
\eqref{eq:rank-one},
\begin{equation}
  \mathbf{b}\,\mathbf{b}^\top
  = \mathbf{I} - \mathbf{R}_{2\times2}^\top \mathbf{R}_{2\times2},
\end{equation}
so everything hinges on how far $\mathbf{R}_{2\times2}$ is from being
orthogonal. At a fronto-parallel view the marker plane is parallel to the image
plane, so the map from the plane to the image is a pure similarity: a rotation
and a uniform scale, with no foreshortening in any direction. Its local
derivative $\mathbf{A}$ is then exactly a scaled rotation, and dividing by its
largest singular value $\gamma$ leaves $\mathbf{R}_{2\times2}$ orthogonal. The
right-hand side vanishes, hence $\mathbf{b}=\mathbf{0}$.

The third column follows: with $\mathbf{b}=\mathbf{0}$ the cross product of the
first two columns is $(0,0,\pm1)^\top$, so $\mathbf{c}=\mathbf{0}$ as well.
Both sign choices produce the \emph{same} matrix, and the solver returns one
pose twice.

Away from $0^\circ$ the plane is foreshortened along one direction only, so
$\mathbf{A}$ is no longer conformal, its two singular values separate, and
$\mathbf{b}$ becomes non-zero. Measuring it along the sweep gives a clean
identity,
\begin{equation}
  \|\mathbf{b}\| = \sin\theta ,
  \label{eq:b-sin}
\end{equation}
which we verified to machine precision at every sampled angle
($\theta=1^\circ$: $0.017452$; $\theta=20^\circ$: $0.342020$;
$\theta=45^\circ$: $0.707107$). So $\|\mathbf{b}\|$ is exactly the
\emph{amount of ambiguity} in the configuration: zero head-on, maximal
edge-on, and growing as the sine of the tilt. The degeneracy at $0^\circ$ is
not a numerical accident but a property of the geometry.

\subsubsection{Why the pose error is exactly twice the viewing angle}

Column three of Table~\ref{tab:sweep} reads $2\theta$ at every sampled angle,
with no rounding error. The reason is that IPPE's two solutions place the
marker normal symmetrically about the line of sight.

Let $\mathbf{d}$ be the unit vector from the marker centre to the camera
centre. By definition of the viewing angle, the true normal $\mathbf{n}_1$
makes an angle $\theta$ with $\mathbf{d}$. Both candidates reproduce the same
local homography, and the only other plane orientation that does so is the one
obtained by tilting the normal by $\theta$ on the \emph{other} side of
$\mathbf{d}$. The two normals are therefore coplanar with $\mathbf{d}$ and
separated by
\begin{equation}
  \angle(\mathbf{n}_1,\mathbf{n}_2) = 2\theta .
\end{equation}

That accounts for the normals, but a rotation has three degrees of freedom and
the normal fixes only two: in principle the two candidates could also differ by
an in-plane spin, which would make the geodesic distance \emph{larger} than
$2\theta$. They do not. Computing the relative rotation
$\mathbf{R}_1^\top\mathbf{R}_2$ along the sweep shows that its axis is, to
machine precision, the marker's $Y$ axis --- the in-plane axis orthogonal to
the tilt direction, the azimuth being $0^\circ$ --- and its angle is exactly
$2\theta$:
\begin{equation}
  \mathbf{R}_1^\top\mathbf{R}_2
  = \exp\!\big([\,2\theta\,\mathbf{u}\,]_\times\big),
  \qquad \mathbf{u}\perp\mathbf{n}_1,\ \mathbf{u}\ \text{in the marker plane}.
\end{equation}
In other words the mirror candidate is the true pose \emph{flipped about an
in-plane axis}, with no extra spin: the whole discrepancy is a re-orientation
of the plane, which is why the geodesic error and the angle between the normals
agree. The consequence is uncomfortable, and is the point of the experiment: at
a modest $\theta=5^\circ$ the wrong candidate is already $10^\circ$ away in 3D,
while its image differs by about $1$ px.

\paragraph{The image is far less sensitive than the pose.} At $5^\circ$ the
two poses are $10^\circ$ apart yet their images differ by about $1$ px. This
gap between 3D separation and image separation is the practical form of the
ambiguity: the pixels simply do not carry enough evidence to decide.

\paragraph{Separability depends on the noise.} With image noise of standard
deviation $\sigma_{px}$, the two candidates are distinguishable only where the
mirror RMSE is comfortably above $\sigma_{px}$. For the $\sigma_{px} = 0.5$ px
used in Experiment 2, that threshold falls at around $2$--$3^\circ$ for this
configuration. Both the distance and the marker size enter here through the
apparent size of the square, so the threshold is a property of the
configuration and not of the method.

\paragraph{The curve saturates.} Beyond $75^\circ$ the mirror RMSE gains only
half a pixel, and by $89^\circ$ the marker is under three pixels wide. Extreme
obliquity separates the candidates no further while making every other aspect
of the problem worse.
\FloatBarrier

% ============================================================
\section{Experiment 2 --- Conditioning and Monte Carlo}
\label{sec:experiment2}
% ============================================================

The second experiment treats pose estimation from the marker as the
least-squares problem of minimizing the reprojection cost $E(\mathbf{R},\mathbf{t})$
of \eqref{eq:cost}. It addresses two questions. First, how stable is that
minimization with respect to noise on the four corners, as measured by the
conditioning of its Jacobian at every viewing angle, and which directions of
the pose are poorly observable? Second, does the uncertainty predicted from the
Jacobian agree with a Monte Carlo simulation in which noise is added to the
corners? Section~\ref{sec:jacobian} derives the Jacobian and
Section~\ref{sec:svd} relates its singular values to the uncertainty of the
estimate. Section~\ref{sec:units} fixes the units in which conditioning is
measured, Section~\ref{sec:frontal} solves the fronto-parallel case in closed
form, and Section~\ref{sec:conditioning} reports the conditioning along the
camera arc. Section~\ref{sec:montecarlo} presents the Monte Carlo validation and
Section~\ref{sec:exp2-discussion} draws the conclusions.

% ------------------------------------------------------------
\subsection{The Residual Jacobian}
\label{sec:jacobian}
% ------------------------------------------------------------

Let $\mathbf{e}\in\R^{8}$ stack the residuals \eqref{eq:residual} of the four
corners, so that $E = \|\mathbf{e}\|^2$. Four corners provide eight equations
for the six unknowns of the pose, so the problem is over-determined by two and
noise cannot in general be fitted exactly.

\paragraph{Local parametrization.} Translation is perturbed additively,
$\mathbf{t}\mapsto\mathbf{t}+\delta\mathbf{t}$. Rotation is perturbed through a
local rotation vector applied on the right,
\begin{equation}
  \mathbf{R}(\delta\boldsymbol{\omega})
  = \mathbf{R}\exp\!\big([\delta\boldsymbol{\omega}]_\times\big),
  \qquad
  \exp\!\big([\delta\boldsymbol{\omega}]_\times\big)
  = \mathbf{I} + [\delta\boldsymbol{\omega}]_\times + O(\|\delta\boldsymbol{\omega}\|^2),
  \label{eq:local-rotation}
\end{equation}
so that $\delta\boldsymbol{\omega}$ is expressed in the marker frame. The six
pose parameters are ordered as
\begin{equation}
  \delta\boldsymbol{\theta} =
  (\delta t_x,\delta t_y,\delta t_z,
   \delta\omega_x,\delta\omega_y,\delta\omega_z)^\top,
  \label{eq:theta}
\end{equation}
with translations in metres and rotations in radians. Euler angles are avoided:
they introduce singularities and make the Jacobian depend on an arbitrary axis
order.

\paragraph{Derivation.} The residual of corner $i$ depends on the pose only
through the camera-frame point $\Pcam^{(i)} = \mathbf{R}\,\Pmark^{(i)} +
\mathbf{t}$, so by the chain rule its Jacobian is the product of two factors.
Substituting \eqref{eq:local-rotation} and using
$[\delta\boldsymbol{\omega}]_\times\Pmark = \delta\boldsymbol{\omega}\times\Pmark
= -[\Pmark]_\times\delta\boldsymbol{\omega}$ gives the first,
\begin{equation}
  \frac{\partial \Pcam^{(i)}}{\partial \delta\boldsymbol{\theta}}
  = \begin{bmatrix}\mathbf{I}_3 & -\mathbf{R}\,[\Pmark^{(i)}]_\times\end{bmatrix}
  \qquad (3\times6).
\end{equation}
Differentiating the pixel projection \eqref{eq:pixel-projection} with respect to
$\Pcam = (X,Y,Z)^\top$ gives the second,
\begin{equation}
  \mathbf{D}_i
  = \begin{bmatrix}
      f_x/Z & 0 & -f_x X/Z^2 \\[2pt]
      0 & f_y/Z & -f_y Y/Z^2
    \end{bmatrix}
  \qquad (2\times3),
\end{equation}
so that
\begin{equation}
  \boxed{\quad
  \mathbf{J}_i = \mathbf{D}_i
  \begin{bmatrix}\mathbf{I}_3 & -\mathbf{R}\,[\Pmark^{(i)}]_\times\end{bmatrix},
  \qquad
  \mathbf{J} = \begin{bmatrix}\mathbf{J}_1\\ \mathbf{J}_2\\ \mathbf{J}_3\\ \mathbf{J}_4\end{bmatrix}
  \in \R^{8\times6}.
  \quad}
  \label{eq:jacobian}
\end{equation}
The first three columns of $\mathbf{J}$ are in pixels per metre and the last
three in pixels per radian. The principal point does not enter $\mathbf{J}$.

\paragraph{Verification.} The implementation (\code{reprojection\_jacobian})
was compared with centred finite differences computed independently of the
project code: the projection was rewritten from scratch and the rotation was
perturbed with SciPy rather than with the project's exponential map. At the
baseline pose and at six viewpoints of the arc, including a camera with
$f_x \neq f_y$ and a shifted principal point, the two agree to a relative
accuracy of $2\cdot10^{-11}$.

% ------------------------------------------------------------
\subsection{Singular Values and Observability}
\label{sec:svd}
% ------------------------------------------------------------

The singular value decomposition
\begin{equation}
  \mathbf{J} = \mathbf{U}\boldsymbol{\Sigma}\mathbf{V}^\top
  = \sum_{k=1}^{6}\sigma_k\,\mathbf{u}_k\mathbf{v}_k^\top,
  \qquad \sigma_1\ge\dots\ge\sigma_6,
\end{equation}
satisfies $\mathbf{J}\mathbf{v}_k = \sigma_k\mathbf{u}_k$. Moving the pose by a
small amount $\varepsilon$ along the right singular vector $\mathbf{v}_k$
changes the residuals by $\varepsilon\,\sigma_k$ in norm. A large $\sigma_k$
therefore marks a direction of the pose that is clearly visible in the image,
and a small one a direction along which the pose can move while the corners
barely do. The smallest singular value $\sigma_{\min} = \sigma_6$ measures how
poorly the worst direction is observed, $\mathbf{v}_6$ identifies that
direction, and the condition number $\kappa = \sigma_1/\sigma_6$ compares the
best and worst observed directions.

\paragraph{From singular values to uncertainty.} Let the observed pixels carry
independent noise $\mathbf{n}$ with $\operatorname{Cov}(\mathbf{n}) =
\sigma_{px}^2\mathbf{I}$. Linearizing the residuals around the true pose,
$\mathbf{e}(\delta\boldsymbol{\theta}) \approx \mathbf{J}\,\delta\boldsymbol{\theta}
- \mathbf{n}$, the minimizer of $E$ is the linear least-squares solution
\begin{equation}
  \delta\boldsymbol{\theta}
  = \big(\mathbf{J}^\top\mathbf{J}\big)^{-1}\mathbf{J}^\top\mathbf{n},
  \label{eq:gn-at-truth}
\end{equation}
whose covariance is
\begin{equation}
  \boxed{\quad
  \operatorname{Cov}(\delta\boldsymbol{\theta})
  = \sigma_{px}^2\big(\mathbf{J}^\top\mathbf{J}\big)^{-1}
  = \sigma_{px}^2\sum_{k=1}^{6}\frac{\mathbf{v}_k\mathbf{v}_k^\top}{\sigma_k^2}.
  \quad}
  \label{eq:covariance}
\end{equation}
The estimate thus scatters along each $\mathbf{v}_k$ with standard deviation
$\sigma_{px}/\sigma_k$: the weakest direction carries the largest error. The
prediction describes the minimizer of $E$; Section~\ref{sec:montecarlo} returns
to this point.

% ------------------------------------------------------------
\subsection{Units and the Dimensionless Jacobian}
\label{sec:units}
% ------------------------------------------------------------

The six parameters do not share a unit, so the singular values of $\mathbf{J}$
mix pixels per metre with pixels per radian, and the condition number depends
on that choice. Measuring translation in millimetres instead of metres moves the
condition number at a $30^\circ$ view from $55$ to $652$, although nothing
physical has changed. More tellingly, a $10$~cm marker at $0.5$~m and a
$20$~cm marker at $1$~m produce exactly the same image, yet in SI units they
have condition numbers $200$ and $100$ head-on.

We therefore measure translation in units of the camera--marker distance
$Z = \|\mathbf{t}\|$ and work with the dimensionless Jacobian
\begin{equation}
  \mathbf{J}_{\mathrm{n}} = \mathbf{J}\,\mathbf{S},
  \qquad
  \mathbf{S} = \operatorname{diag}(Z,Z,Z,1,1,1).
  \label{eq:dimensionless}
\end{equation}
A translation of one unit then displaces the marker by about one radian as seen
from the camera, so all six parameters are dimensionless and comparable. The
singular values of $\mathbf{J}_{\mathrm{n}}$ are in pixels and its condition
number $\kappa$ is a pure number: identical images give identical values, and
$\kappa$ does not depend on the focal length. The rotation columns are
unaffected. From here on, $\sigma_k$, $\mathbf{v}_k$ and $\kappa$ refer to
$\mathbf{J}_{\mathrm{n}}$; the covariances of \eqref{eq:covariance} remain in
metres and radians.

% ------------------------------------------------------------
\subsection{The Fronto-Parallel Case in Closed Form}
\label{sec:frontal}
% ------------------------------------------------------------

At a fronto-parallel view the Jacobian can be analysed by hand. Let
$f_x = f_y = f$, let the marker face the camera at distance $Z$ on the optical
axis, $\mathbf{R} = \operatorname{diag}(1,-1,-1)$ and $\mathbf{t} =
(0,0,Z)^\top$, and let the corners be $\Pmark = (x,y,0)^\top$ with
$x,y\in\{\pm h\}$ and $h = L/2$. This is the $0^\circ$ pose of the camera arc.
All four corners lie at depth $Z$, and \eqref{eq:jacobian} gives, for each
corner and in the order \eqref{eq:theta}, the two rows
\begin{align}
  \text{$u$-row:}\quad
  &\frac{f}{Z}\Big(1,\ 0,\ -\frac{x}{Z},\ \frac{xy}{Z},\ -\frac{x^2}{Z},\ -y\Big),
  \nonumber\\
  \text{$v$-row:}\quad
  &\frac{f}{Z}\Big(0,\ 1,\ \frac{y}{Z},\ -\frac{y^2}{Z},\ \frac{xy}{Z},\ -x\Big).
  \label{eq:frontal-rows}
\end{align}
By the symmetry of the square, sums of odd powers of $x$ and $y$ vanish, and
$\mathbf{J}^\top\mathbf{J}$ becomes block diagonal:
\begin{equation}
  \frac{\mathbf{J}^\top\mathbf{J}}{(f/Z)^2}
  =
  \underbrace{\begin{bmatrix} 4 & -4s \\ -4s & 8s^2 \end{bmatrix}}_{(t_x,\,\omega_y)}
  \oplus
  \underbrace{\begin{bmatrix} 4 & -4s \\ -4s & 8s^2 \end{bmatrix}}_{(t_y,\,\omega_x)}
  \oplus
  \underbrace{\Big[\tfrac{8h^2}{Z^2}\Big]}_{t_z}
  \oplus
  \underbrace{\big[8h^2\big]}_{\omega_z},
  \qquad s = \frac{h^2}{Z}.
\end{equation}
Each out-of-plane rotation is coupled to one sideways translation, while depth
and in-plane spin decouple. Since the marker is small compared with its
distance, $s \ll 1$, and to leading order the singular values of the
dimensionless Jacobian \eqref{eq:dimensionless} are
\begin{equation}
  \sigma_{t_x} = \sigma_{t_y} \approx 2f,
  \qquad
  \sigma_{t_z} = \sigma_{\omega_z} = \frac{2\sqrt2\,fh}{Z},
  \qquad
  \sigma_{\min} \approx \frac{2fh^2}{Z^2},
  \qquad
  \boxed{\;\kappa \approx \Big(\frac{Z}{h}\Big)^2\;}
  \label{eq:frontal-sigmas}
\end{equation}
The condition number head-on depends only on the marker's apparent half-size
$h/Z$: \emph{the conditioning degrades as the inverse square of how large the
marker appears.}

\paragraph{Why the tilts are the weak directions.} The eigenvectors of the small
eigenvalues are the two out-of-plane tilts, each accompanied by a small sideways
translation. Rows \eqref{eq:frontal-rows} show why. Since $x^2 = h^2$ at every
corner, a tilt $\delta\omega_y$ moves all four corners horizontally by the same
amount, $-fh^2/Z^2$ per radian; that uniform shift is reproduced exactly by a
sideways translation and therefore cancels. What remains is the vertical motion,
whose sign follows $xy$: two opposite corners move up and the other two down, a
small keystone distortion of the square. Head-on, that keystone is the only
first-order evidence of a tilt, and its size is of order $h/Z$ smaller than the
effect of translating the marker.

Table~\ref{tab:frontal} compares \eqref{eq:frontal-sigmas} with the
implementation for $f = 800$~px, $h = 0.05$~m and $Z = 0.5$~m. The small
differences are the neglected terms of order $s^2$.

\begin{table}[h!]
\centering\small
\begin{tabularx}{\linewidth}{@{}Yrr@{}}
\toprule
\rowcolor{accent}\thcell{Quantity}&\thcell{Closed form}&\thcell{Implementation}\\
\midrule
$\sigma_{t_x} = \sigma_{t_y}$ & $1600.0$ & $1600.08$\\
\rowcolor{lightbg} $\sigma_{t_z}$ & $226.274$ & $226.274$\\
$\sigma_{\omega_z}$ & $226.274$ & $226.274$\\
\rowcolor{lightbg} $\sigma_{\min}$ (twice) & $16.000$ & $15.9992$\\
$\kappa$ & $(Z/h)^2 = 100$ & $100.010$\\
\bottomrule
\end{tabularx}
\caption{Fronto-parallel view: closed-form singular values of the dimensionless
Jacobian against the implementation.}
\label{tab:frontal}
\end{table}

% ------------------------------------------------------------
\subsection{Conditioning Along the Camera Arc}
\label{sec:conditioning}
% ------------------------------------------------------------

Away from the fronto-parallel view the algebra no longer closes, and the
Jacobian is evaluated numerically (\code{run\_conditioning\_sweep}) at the true
pose of every viewpoint of the arc of Experiment~1: radius $0.5$~m, azimuth
$0^\circ$, one viewpoint per degree from $0^\circ$ to $89^\circ$.
Figure~\ref{fig:conditioning} shows the result and Table~\ref{tab:conditioning}
lists representative values.

\resultfigure{exp2_conditioning.png}
 {Singular values, condition number and composition of the weakest direction against the viewing angle.}
 {Conditioning of the dimensionless reprojection Jacobian along the camera arc.
 Top: the six singular values; $\sigma_1 = \sigma_2$ and $\sigma_3 \approx \sigma_4$,
 so $\sigma_2$ is drawn $10\%$ lower and $\sigma_3$ $10\%$ higher for visibility.
 Middle: the condition number $\kappa = \sigma_1/\sigma_6$. Bottom: the share of
 each pose parameter in the least observable direction $\mathbf{v}_6$ (squared
 components); $t_x$, $t_y$, $t_z$ and $\omega_y$ together stay below $0.01\%$ at
 every angle. Generated by \code{experiment2.py}.}
 {fig:conditioning}

\begin{table}[h!]
\centering\small
\begin{tabularx}{\linewidth}{@{}rrrrrrrrY@{}}
\toprule
\rowcolor{accent}\thcell{$\theta$}&\thcell{$\sigma_1$}&\thcell{$\sigma_2$}&\thcell{$\sigma_3$}&\thcell{$\sigma_4$}&\thcell{$\sigma_5$}&\thcell{$\sigma_6$}&\thcell{$\kappa$}&\thcell{Weakest direction}\\
\midrule
$0^\circ$ & 1600.1 & 1600.1 & 226.3 & 226.3 & 16.0 & 16.0 & 100.0 & $\omega_x + 0.010\,t_y$\\
\rowcolor{lightbg} $5^\circ$ & 1600.3 & 1600.3 & 226.1 & 226.1 & 18.9 & 18.5 & 86.4 & $\omega_x - 0.04\,\omega_z$\\
$15^\circ$ & 1601.7 & 1601.7 & 224.6 & 224.5 & 34.1 & 32.3 & 49.6 & $\omega_x - 0.13\,\omega_z$\\
\rowcolor{lightbg} $30^\circ$ & 1606.2 & 1606.1 & 219.5 & 219.1 & 61.8 & 58.2 & 27.6 & $\omega_x - 0.26\,\omega_z$\\
$45^\circ$ & 1612.3 & 1612.1 & 210.8 & 210.0 & 89.7 & 85.1 & 18.9 & $\omega_x - 0.38\,\omega_z$\\
\rowcolor{lightbg} $60^\circ$ & 1618.4 & 1618.2 & 198.7 & 197.0 & 116.4 & 111.6 & 14.5 & $\omega_x - 0.51\,\omega_z$\\
$75^\circ$ & 1622.9 & 1622.7 & 183.2 & 180.3 & 141.0 & 136.8 & 11.9 & $\omega_x - 0.65\,\omega_z$\\
\rowcolor{lightbg} $89^\circ$ & 1624.5 & 1624.3 & 168.3 & 162.8 & 159.0 & 157.2 & 10.3 & $\omega_z - 0.28\,\omega_x$\\
\bottomrule
\end{tabularx}
\caption{Singular values (pixels) and condition number of the dimensionless
Jacobian along the arc, with the dominant components of the least observable
direction.}
\label{tab:conditioning}
\end{table}

\paragraph{Which motion each singular value belongs to.} The singular vectors
attach every singular value to one motion of the marker, and the geometry
explains why they come in pairs.
\begin{itemize}[leftmargin=*]
  \item $\sigma_1,\sigma_2$ belong to the two sideways translations. Sliding the
        marker moves all four corners together; the two are equal because a
        square is unchanged by a $90^\circ$ turn, and they stay nearly constant
        because they depend only on $f$ and $Z$, both fixed along the arc.
  \item $\sigma_3,\sigma_4$ belong to depth and in-plane spin. Bringing the
        marker closer by one unit of $Z$ zooms the image and moves each corner,
        at distance $r = h\sqrt2$ from the centre, outwards by $fr/Z$ pixels;
        spinning it by one radian moves each corner along a circle by the same
        amount. Both decrease as the marker is foreshortened.
  \item $\sigma_5,\sigma_6$ belong to the two out-of-plane tilts. They are equal
        head-on by the same symmetry and small because a tilt then shows only as
        a keystone. The arc tilts the marker about its own $Y$ axis, which breaks
        the symmetry and separates them.
\end{itemize}

\paragraph{Reading the condition number.} A larger $\kappa$ means a worse
conditioned problem: the pose is then determined very unevenly, precisely in
some directions and loosely in others. The fronto-parallel view is the worst
case, $\kappa = 100$, and the conditioning improves steeply with obliquity:
$\kappa$ halves by $15^\circ$ and reaches about $28$ at $30^\circ$, after which
the gains are small, down to $10$ at $89^\circ$. Since $\sigma_1$ is essentially
constant, the whole trend is carried by $\sigma_6$, which grows tenfold from
$16$ to $157$. The reason is geometric. Head-on, a small out-of-plane rotation
moves the corners only along the line of sight and is visible solely through the
keystone. Once the plane is oblique, the same rotation also moves the corners
perpendicular to the line of sight, with opposite signs on opposite edges, which
no translation can reproduce, so the evidence becomes first order in $h/Z$.

\paragraph{The least observable direction.} It is always a rotation and never a
translation: the three translations and $\omega_y$ together never exceed
$0.01\%$ of $\mathbf{v}_6$. Up to about $20^\circ$ it is a pure out-of-plane tilt
$\omega_x$, the marker nodding about its horizontal axis, perpendicular to the
axis about which the camera moves along the arc. As the view becomes oblique it
acquires an in-plane spin component, and near $89^\circ$ it becomes mostly
$\omega_z$: seen almost edge-on, spinning the marker in its own plane moves the
corners along the line of sight, which the image hardly registers.

% ------------------------------------------------------------
\subsection{Monte Carlo Validation}
\label{sec:montecarlo}
% ------------------------------------------------------------

\subsubsection{The estimator}

Equation \eqref{eq:covariance} describes the \emph{minimizer of $E$}, so the
simulation must produce exactly that estimate. Since $E$ is a sum of squared
non-linear residuals, it is minimized with Gauss--Newton, which requires only
the Jacobian \eqref{eq:jacobian}. Starting from a pose, each iteration evaluates
$\mathbf{e}$ and $\mathbf{J}$, computes the step
\begin{equation}
  \delta\boldsymbol{\theta}
  = -\big(\mathbf{J}^\top\mathbf{J}\big)^{-1}\mathbf{J}^\top\mathbf{e},
  \label{eq:gauss-newton}
\end{equation}
which minimizes the linearized cost $\|\mathbf{e}+\mathbf{J}\,\delta\boldsymbol{\theta}\|^2$,
and updates $\mathbf{t}\leftarrow\mathbf{t}+\delta\mathbf{t}$ and
$\mathbf{R}\leftarrow\mathbf{R}\exp([\delta\boldsymbol{\omega}]_\times)$, until
$\|\delta\boldsymbol{\theta}\| < 10^{-10}$. At the true pose, where
$\mathbf{e} = -\mathbf{n}$, the step \eqref{eq:gauss-newton} is exactly
\eqref{eq:gn-at-truth}: the prediction is one Gauss--Newton step, and the
simulation tests whether the full, iterated minimization behaves as that
first-order picture states.

Gauss--Newton is local, and for a planar marker $E$ has two valleys, around the
true pose and around its mirror. IPPE supplies one starting point in each: both
candidates are refined and the one with the lower cost is kept
(\code{estimate\_pose}). The result does not depend on the starting point within
a valley: on the same noisy data, Gauss--Newton started from the IPPE candidates
and from the true pose reached the same estimate in all of $300$ trials at
$5^\circ$ and at $30^\circ$, to within $10^{-11}$ degrees. Raw IPPE, by contrast,
lies $2.7^\circ$ from that minimizer at $5^\circ$ (median), more than the
predicted scatter itself, so it could not be used in place of the minimizer.

\subsubsection{Procedure}

At each viewing angle $N = 2500$ trials are run. Each trial adds independent
Gaussian noise with $\sigma_{px} = 0.5$~px to the eight corner coordinates,
estimates the pose as above, and records its error in the coordinates of the
Jacobian, $\delta\mathbf{t} = \hat{\mathbf{t}} - \mathbf{t}$ and
$\delta\boldsymbol{\omega} = \log(\mathbf{R}^\top\hat{\mathbf{R}})$. A trial is
counted as a \emph{mirror flip} when the estimated rotation is closer to the
mirror pose (IPPE's second candidate on exact data) than to the true one; at
$0^\circ$ the two coincide and no flip is possible. Statistics are computed on
the trials that remained in the valley of the true pose, since
\eqref{eq:covariance} describes a single minimum, and flips are reported
separately. Three comparisons are made: the ratio of measured to predicted
standard deviation for each parameter, whose statistical uncertainty over $N$
trials gives a $\pm2\sigma$ band of $\pm2/\sqrt{2(N-1)} = \pm2.8\%$; the measured
scatter along $\mathbf{v}_6$ against $\sigma_{px}/\sigma_6$; and the fraction of
flips. The viewing angles are dense near head-on, $0^\circ$, $0.5^\circ$,
$1^\circ$, $2^\circ$, \dots, $8^\circ$, where the mirror can win, and every
$5^\circ$ from $10^\circ$ to $85^\circ$.

\subsubsection{Results}

\resultfigure{exp2_montecarlo.png}
 {Measured against predicted scatter per parameter, scatter along the weakest direction, and mirror flips against the viewing angle.}
 {Monte Carlo validation, $2500$ trials per angle, $\sigma_{px} = 0.5$~px.
 Top: ratio of measured to predicted standard deviation for each pose
 parameter; the grey band is the $\pm2\sigma$ statistical uncertainty. Middle:
 standard deviation along the least observable direction, measured against the
 prediction $\sigma_{px}/\sigma_6$. Bottom: percentage of trials in which the
 mirror pose was selected, on $0^\circ$--$15^\circ$. Generated by
 \code{experiment2.py}.}
 {fig:montecarlo}

\begin{table}[h!]
\centering\small
\begin{tabularx}{\linewidth}{@{}rYrrrr@{}}
\toprule
\rowcolor{accent}\thcell{$\theta$}&\thcell{Mirror flips}&\thcell{Ratio $\omega_y$}&\thcell{Ratio $t_x$}&\thcell{$\mathbf{v}_6$ pred.\ ($^\circ$)}&\thcell{$\mathbf{v}_6$ meas.\ ($^\circ$)}\\
\midrule
$0^\circ$ & $0$ & $1.023$ & $1.012$ & $1.791$ & $1.824$\\
\rowcolor{lightbg} $0.5^\circ$ & $953$ ($38\%$) & $0.679$ & $0.841$ & $1.788$ & $1.881$\\
$1^\circ$ & $694$ ($28\%$) & $0.694$ & $0.841$ & $1.778$ & $1.801$\\
\rowcolor{lightbg} $2^\circ$ & $333$ ($13\%$) & $0.817$ & $0.901$ & $1.744$ & $1.849$\\
$3^\circ$ & $127$ ($5\%$) & $0.903$ & $0.952$ & $1.690$ & $1.758$\\
\rowcolor{lightbg} $4^\circ$ & $35$ ($1.4\%$) & $0.996$ & $1.006$ & $1.622$ & $1.608$\\
$6^\circ$ & $6$ ($0.2\%$) & $1.017$ & $0.994$ & $1.466$ & $1.507$\\
\rowcolor{lightbg} $10^\circ$ & $0$ & $1.021$ & $0.995$ & $1.164$ & $1.168$\\
$30^\circ$ & $0$ & $1.005$ & $0.996$ & $0.492$ & $0.495$\\
\rowcolor{lightbg} $60^\circ$ & $0$ & $1.001$ & $0.996$ & $0.257$ & $0.258$\\
$85^\circ$ & $0$ & $0.997$ & $1.000$ & $0.188$ & $0.186$\\
\bottomrule
\end{tabularx}
\caption{Monte Carlo results at representative angles. Ratios are measured over
predicted standard deviations. The other four parameters stay close to $1$ at
every angle; from $4^\circ$ onwards all six lie within $3.6\%$ of $1$.}
\label{tab:montecarlo}
\end{table}

\subsubsection{Interpretation}

\paragraph{The prediction holds.} At $0^\circ$ and from $4^\circ$ to $85^\circ$,
the measured scatter agrees with $\sigma_{px}^2(\mathbf{J}^\top\mathbf{J})^{-1}$
within statistical noise, both parameter by parameter and along the least
observable direction, whose standard deviation falls from $1.8^\circ$ head-on to
$0.19^\circ$ at $85^\circ$. The error also concentrates where the conditioning
analysis placed it: in the out-of-plane rotations, while the sideways
translations scatter by about $0.2$~mm at every angle.

\paragraph{Where it stops holding.} Between $0.5^\circ$ and $3^\circ$ three
effects appear together, with a single cause: the mirror pose lies within the
noise.
\begin{itemize}[leftmargin=*]
  \item \emph{Mirror flips.} The mirror is the true pose rotated by $2\theta$
        about the marker's $Y$ axis. Near head-on this is within the tilt
        scatter of about $1.8^\circ$, and noise frequently makes the mirror the
        better fit: $38\%$ of the trials at $0.5^\circ$, tending to one half as
        $\theta\to0$, and none from $8^\circ$ onwards. In those trials the noisy
        corners genuinely fit the mirror better, so no estimator based on the
        four corners alone could avoid the error. This is the ambiguity of
        Experiment~1 measured under noise.
  \item \emph{Truncated statistics.} Because the mirror lies along $\omega_y$,
        the trials that flip are those with the largest $\omega_y$ errors, and
        excluding them makes the remaining distribution too narrow: ratios down
        to $0.68$. The translation $t_x$, coupled to $\omega_y$
        (Section~\ref{sec:frontal}), follows. A ratio below one therefore
        signals that the single-minimum prediction no longer applies, not that
        $\omega_y$ is better observed.
  \item \emph{A slightly larger scatter along $\mathbf{v}_6$}, by $3$--$6\%$:
        with two minima this close, $E$ is flatter than the quadratic bowl
        assumed by the first-order analysis.
\end{itemize}
At exactly $0^\circ$ the two valleys merge into one, and the prediction holds
again.

% ------------------------------------------------------------
\subsection{Discussion}
\label{sec:exp2-discussion}
% ------------------------------------------------------------

\paragraph{Fronto-parallel is the worst case.} Contrary to the intuition that
looking straight at a marker is the easy configuration, the fronto-parallel view
is the least well conditioned. Its condition number is $(Z/h)^2$: doubling the
distance, or halving the marker, makes it four times worse.

\paragraph{What is lost is orientation, not position.} The least observable
direction is always a rotation. With $\sigma_{px} = 0.5$~px, the out-of-plane
tilt is uncertain by $1.8^\circ$ head-on, $0.5^\circ$ at $30^\circ$ and
$0.26^\circ$ at $60^\circ$, while the sideways position is determined to a
fraction of a millimetre throughout. A modest obliquity of $20^\circ$--$30^\circ$
secures most of the improvement.

\paragraph{The first-order theory is validated.} With the pose estimated as the
minimizer of $E$, the Jacobian predicts the Monte Carlo scatter within
statistical noise everywhere except in a narrow band near head-on.

\paragraph{The two experiments describe one phenomenon.} Experiment~1 showed that
the two IPPE candidates differ by the sign of a vector of length
$\sin\theta$, \eqref{eq:b-sin}, so that near head-on the image
reveals how much the marker is tilted but hardly towards which side.
Experiment~2 shows that the same out-of-plane tilt is the least observable
direction of the least-squares problem, with a first-order signature of relative
size $h/Z$ only. The Monte Carlo joins the two: the prediction fails exactly in
the band where the mirror solution lies within the noise and is selected in up
to $38\%$ of the trials.
\FloatBarrier

\end{document}
